\section{Metodologia}
\label{sec:metodologia}
O protocolo separa a fixação dos hiperparâmetros, o ajuste no treino, a calibração em junho e a avaliação nos demais recortes. AE e IF utilizam a mesma representação e pontuam toda a população de cada janela.

\subsection{Dados e recortes temporais}
Utilizamos os Parquets diários de transações da Ethereum Mainnet disponibilizados pelo AWS Public Blockchain Data \cite{AWSBlockchain}. Os 365 dias de 2025 totalizam 523.937.952 registros na fonte utilizada. O processamento interpreta os timestamps em UTC e inclui todas as transações disponíveis de cada dia, preservando o contexto intrabloco usado na construção dos atributos.

\begin{table}[htbp]
\centering\small
\caption{Recortes de 2025, datas inclusivas em UTC.}
\label{tab:recortes}
\begin{tabular}{@{}llrl@{}}
\toprule
Recorte & Período & Transações & Uso \\
\midrule
Treino & 01/01--30/04 & 147.512.905 & Ajuste \\
Pectra & 01/05--31/05 & 41.982.829 & Retrospectiva \\
Validação & 01/06--30/06 & 41.950.663 & Calibração \\
Teste & 01/07--30/11 & 241.020.375 & Avaliação \\
Fusaka/BPO1 & 01/12--09/12 & 14.264.473 & Transição \\
Estresse & 10/12--31/12 & 37.206.707 & Estresse temporal \\
\bottomrule
\end{tabular}
\end{table}

A Tabela~\ref{tab:recortes} distingue o período de ajuste, janeiro a abril, da calibração em junho e do teste de julho a novembro. Maio é analisado retrospectivamente com os cortes calibrados em junho. Dezembro é dividido em uma janela de transição, de 1 a 9, e uma janela de estresse, de 10 a 31. Os marcos Pectra, ativado em 7 de maio, e Fusaka/BPO1, em 3 e 9 de dezembro, contextualizam a organização temporal \cite{Pectra2025,Fusaka2025}. As comparações caracterizam os regimes observados, sem atribuir a uma atualização isolada as mudanças encontradas.

Cada registro recebe um identificador determinístico $r=b\,2^{32}+i$, em que $b$ é o número do bloco e $i$ é o índice intrabloco. A auditoria verifica limites inteiros, consistência temporal, esquema, chaves, nulos e finitude dos valores. Os manifestos vinculam os arquivos às suas contagens e hashes, permitindo conferir a população processada em cada etapa.

\subsection{Representação e pré-processamento}
Cada transação é representada pelos 21 atributos da configuração B, agrupados na Tabela~\ref{tab:atributos}. Os atributos combinam execução e valor, condições de taxa, tipos de transação, contexto do bloco e componentes cíclicos do tempo. Contagens e ranks intrabloco são construídos com todos os registros disponíveis do bloco. Trata-se de uma representação posterior à inclusão e execução da transação.

\begin{table}[htbp]
\centering\small
\caption{Atributos utilizados pelos modelos. $\ell(x)=\ln(1+x)$.}
\label{tab:atributos}
\begin{tabularx}{\linewidth}{@{}p{2.7cm}X@{}}
\toprule
Grupo & Atributos \\
\midrule
Execução e valor & $\ell$(gás usado); gás usado/limite; $\ell$(valor); $\ell$(nonce); status do recibo. \\
Taxas & $\ell$(preço efetivo); $\ell$(taxa máxima); $\ell$(prioridade máxima); preço efetivo/taxa máxima. \\
Tipos & Indicadores dos tipos 0, 1 e 3. \\
Contexto do bloco & Posição relativa; $\ell$(número de transações); $\ell$(contagens por remetente e destinatário); rank percentual do preço efetivo. \\
Tempo & Seno/cosseno da hora e do dia da semana. \\
\bottomrule
\end{tabularx}
\end{table}


Os nove atributos logarítmicos utilizam $\ell(x)=\ln(1+x)$ e são normalizados por $z_j(x)=(\ell(x_j)-m_j)/\mathrm{IQR}_j$, em que $m_j$ e $\mathrm{IQR}_j$ são a mediana e o intervalo interquartil estimados exclusivamente no treino. Um IQR nulo recebe escala unitária. Para taxa máxima e prioridade máxima, tipos legados 0/1 são excluídos do ajuste da escala e mantidos em zero na transformação. Indicadores, razões, ranks e componentes cíclicos preservam suas escalas. Os tipos 0, 1 e 3 possuem indicadores explícitos; tipos 2 e 4 são representados pelos demais atributos. As matrizes são armazenadas em precisão dupla e convertidas para \texttt{float32} na entrada dos modelos.

\subsection{Origem e fixação dos hiperparâmetros}
\label{sec:hiperparametros}
Os hiperparâmetros são pré selecionados em um experimento anterior da dissertação e mantidos fixos neste estudo. Naquele experimento, a grade do AE combinou três arquiteturas simétricas, duas taxas de aprendizado e dois valores de \textit{weight decay}; a grade do IF combinou três quantidades de árvores e quatro tamanhos de amostra por árvore. Cada algoritmo teve 12 candidatos, ajustados em cinco representações transacionais, totalizando 120 ajustes de busca. A escolha foi realizada por algoritmo, pela maior média de AP na validação semântica das cinco representações, com desempate determinístico.

Transferimos os vencedores globais dessa seleção (Tabela~\ref{tab:hiperparametros}) para um novo ajuste de um AE e um IF na configuração B de 2025. A escolha precede este experimento. O Label Cloud é aplicado exclusivamente à avaliação posterior, preservando sua independência em relação aos parâmetros, ao treinamento e aos limiares.

\begin{table}[htbp]
\centering\small
\caption{Hiperparâmetros preselecionados e regras de ajuste do artigo.}
\label{tab:hiperparametros}
\begin{tabularx}{\linewidth}{@{}>{\raggedright\arraybackslash}p{1.3cm}>{\raggedright\arraybackslash}p{4.1cm}>{\raggedright\arraybackslash}X@{}}
\toprule
Modelo & Parâmetro ou regra & Valor \\
\midrule
AE & Arquitetura & $21\to16\to8\to4\to8\to16\to21$ \\
AE & Learning rate; weight decay & $3\times10^{-4}$; zero \\
AE & Otimizador; perda & AdamW; MSE de reconstrução \\
AE & Batch; limite de épocas & 65.536; 200 \\
AE & Patience; melhora mínima & 12 épocas; $10^{-5}$ \\
IF & Árvores; amostras por árvore & 100; até 2.048 \\
IF & Alvo da população de ajuste & 500.000 transações \\
Ambos & Semente; representação & 42; B com 21 atributos \\
\bottomrule
\end{tabularx}
\end{table}


\subsection{Ajuste dos modelos e produção dos escores}
O AE possui arquitetura $21\!\to\!16\!\to\!8\!\to\!4\!\to\!8\!\to\!16\!\to\!21$, LeakyReLU com inclinação 0,1 nas camadas intermediárias e saída linear. O ajuste emprega AdamW, learning rate $3\times10^{-4}$, weight decay zero, lotes de 65.536 e precisão mista na GPU. A reserva interna utiliza $r\bmod10=0$ e permanece dentro do intervalo de treino; os demais registros participam do ajuste dos pesos. O limite é de 200 épocas, com \textit{early stopping} após 12 épocas sem melhora aceita superior a $10^{-5}$ na perda interna. O estado aceito por esse critério é utilizado na inferência.

O escore do AE é o erro quadrático médio de reconstrução dos 21 atributos:
\begin{equation}
 s_{AE}(x)=\frac{1}{21}\sum_{j=1}^{21}\bigl(x_j-\widehat{x}_j\bigr)^2.
\end{equation}
O processamento por lotes utiliza matrizes mapeadas em disco; checkpoints preservam o modelo, o otimizador e os estados necessários à retomada.

O IF utiliza 100 árvores, até 2.048 exemplos por árvore, todos os atributos, \texttt{bootstrap=False} e \texttt{contamination=auto}. Seu conjunto de ajuste é obtido pela regra determinística $h(r)\bmod N<\min(N,500.000)$, em que $N$ é a quantidade de transações do treino e $h$ é a função de hash utilizada pelo motor de dados. A quantidade efetivamente selecionada é registrada no manifesto e pode diferir do alvo. Essa seleção define a população disponível para ajuste; as 2.048 amostras por árvore são uma decisão distinta. O escore é o negativo de \texttt{score\_samples}. Ambos os algoritmos usam semente 42 e pontuam todos os registros das janelas, com escores maiores indicando maior anormalidade.

\subsection{Calibração e concordância das sinalizações}
Para cada algoritmo $m$, os cortes $\tau_{m,q}$ são os quantis exatos Q99, Q99,5 e Q99,9 de seus escores na validação de junho. A regra de sinalização é $s_m(x)\geq\tau_{m,q}$. Os cortes permanecem congelados nas demais janelas, preservando também os empates no limiar. Q99,5 é o ponto de operação principal; os demais cortes permitem examinar a sensibilidade ao volume de investigação.

Para um recorte $t$ e um quantil $q$, sejam $A_{t,q}$ e $I_{t,q}$ os conjuntos sinalizados por AE e IF. Registramos suas taxas, interseção, união e parcelas exclusivas. A concordância é $J(A_{t,q},I_{t,q})=|A_{t,q}\cap I_{t,q}|/|A_{t,q}\cup I_{t,q}|$, com valor 1 para dois conjuntos vazios. A interseção prioriza concordância; a união permite investigar qualquer sinalização.

\subsection{Referência auxiliar e avaliação da ordenação}
O CSV histórico do Etherscan Label Cloud \cite{EtherscanLabels} contém 3.913 linhas e 3.750 endereços únicos. Os endereços são normalizados para minúsculas e suas ocorrências agrupadas antes do vínculo pelo remetente. Uma transação é associada quando seu remetente possui a categoria \texttt{mev\_bot}; as demais são denominadas não associadas. A referência caracteriza uma associação histórica de endereços, não a confirmação de um ataque individual. Sua identificação por SHA-256 fixa o arquivo utilizado.

Seja $M_t$ o conjunto de associadas. Calculamos $A_{t,q}\cap M_t$, $I_{t,q}\cap M_t$, a interseção $(A_{t,q}\cap I_{t,q})\cap M_t$ e a união $(A_{t,q}\cup I_{t,q})\cap M_t$, além das parcelas exclusivas. O gráfico de composição normaliza a união associada para 100\% e informa seu tamanho absoluto.

ROC-AUC e AP são calculadas a partir de todos os grupos de escores empatados. A redução de pontos é restrita à apresentação gráfica das curvas ROC e Precision-Recall. Para cada limiar, a \textit{precision} auxiliar é a proporção de associados entre os sinalizados, e o \textit{recall} é a fração de associados recuperada. A AP integra os incrementos de recall ponderados pela precision correspondente. Também são registrados a área PR trapezoidal e os resultados sob orçamentos fixos no topo do ranking. O conjunto dos 1\% maiores escores utiliza $k=\lceil0{,}01N_t\rceil$, com desempate por $r$; seu \textit{lift} compara a concentração de associados à prevalência no recorte. Essas medidas distinguem a qualidade da ordenação e a composição dos conjuntos priorizados \cite{Saito2015}.

\subsection{Caracterização das mudanças de distribuição}
Comparamos atributos e escores ao treino por PSI, com intervalos definidos pelos decis do treino e piso de probabilidade $10^{-6}$, e distância de Kolmogorov--Smirnov aproximada. Esta última avalia a maior diferença absoluta entre as distribuições acumuladas em até 101 quantis do treino. As fronteiras permanecem fixas, caracterizando mudanças marginais entre recortes.

\subsection{Reprodutibilidade}
O código está disponível no \href{https://github.com/marcelocorni/anomalias-ae-if-2025}{repositório público no GitHub}, com instruções de execução. O arquivo \texttt{uv.lock} fixa as dependências, e as configurações registram atributos, recortes, semente e hiperparâmetros. Os vencedores da seleção anterior e o CSV do Label Cloud são identificados por SHA-256. O perfil da revisão integra processamento, avaliações, interseções e figuras; manifestos registram hashes, contagens e tempos por etapa, permitindo retomadas conferidas. Os Parquets da AWS e o CSV histórico são entradas externas, com preparação descrita na documentação.
